% P15-U1: Durrett 2.1.21 / Appendix A.3 and the complete 2.1.22 coding construction.
\chapter{Countable Product Measures and Measurable Real Codings}
\label{sec:p15-countable-product-measures}

Consistent probabilities for finite observations must still be shown to
extend to one countably additive probability on infinite sequences. This
chapter proves Durrett's real-coordinate Kolmogorov extension theorem,
including its compact approximation and diagonal argument
\cite[thm.~2.1.21, p.~53; app.~A.3, pp.~464--466]{DurrettPTEAuthor2019}.
It then completes the distance-coordinate and binary-coding construction
of Theorem 2.1.22, explicitly identifying its Borel image and measurable
inverse \cite[p.~54]{DurrettPTEAuthor2019}.

The index set here is \(\mathbb N_+=\{1,2,\ldots\}\). This is a
countable-index theorem, with a specified product sigma-algebra. The
construction uses the earlier measure-extension theorem; finite product
measures alone are not treated as an infinite-product existence proof.

\section{Consistent marginals and the two cylinder algebras}
\label{sec:p15-cylinder-algebras}

\begin{hbdefinition}{def:p15-countable-cylinders}{Real sequences, rectangles and Borel cylinders}
Let \(\Omega_\infty=\mathbb R^{\mathbb N_+}\), the set of all
real sequences. Write \(\pi_n(\omega)=(\omega_1,\ldots,\omega_n)\)
for its first-\(n\) projection, and \(X_i(\omega)=\omega_i\)
for its \(i\)-th coordinate. A \emph{Borel cylinder} is a set
\(\pi_n^{-1}(B)\) with \(n\geq1\), \(B\in\mathcal B(\mathbb R^n)\).
Let \(\mathcal C\) be the collection of all these cylinders.

Use \(I(a,b)=\{x\in\mathbb R:a<x\leq b\}\) for
\(-\infty\leq a<b\leq\infty\), with infinite endpoints interpreted
as in Definition~\ref{def:p15-stieltjes-intervals}.
Let \(\mathcal S\) consist of the empty set and the rectangular cylinders
\[
 \pi_n^{-1}\left(\prod_{i=1}^n I(a_i,b_i)\right).
\]
The \emph{rectangular cylinder algebra} \(\mathcal A\) consists of
finite disjoint unions of members of \(\mathcal S\). It is the
source's initial algebra. Define the product sigma-algebra by
\(\mathcal F_\infty=\sigma(\mathcal S)\).
The larger Borel-cylinder algebra \(\mathcal C\) is also used below,
so compact closed bases can be assigned a coherent finite-dimensional
mass before a measure on \(\mathcal F_\infty\) has been constructed.
Neither algebra is asserted to contain every subset of \(\Omega_\infty\).
\index{cylinder set}\index{product sigma-algebra!countable}
\end{hbdefinition}

\begin{hbdefinition}{def:p15-consistent-marginals}{Consistent finite-dimensional probability laws}
For each \(n\geq1\), let \(\mu_n\) be a probability measure
on \(\mathcal B(\mathbb R^n)\). The source's consistency condition is
\[
 \mu_{n+1}\left(\prod_{i=1}^n I(a_i,b_i)\times\mathbb R\right)
            =\mu_n\left(\prod_{i=1}^n I(a_i,b_i)\right)
\]
for every \(n\) and every such rectangle. Testing bounded real
endpoints is already sufficient, as proved next. No independence is
assumed: each \(\mu_n\) may describe dependent coordinates.
\index{consistent finite-dimensional distributions}
\end{hbdefinition}

\begin{hblemma}{thm:p15-cylinder-consistency}{proved=annex:p15.cylinder-consistency}{Coherent cylinder masses and finite additivity}
The family \(\mathcal S\) is a semialgebra, \(\mathcal A\) and
\(\mathcal C\) are algebras, and
\[
 \mathcal A\subseteq\mathcal C\subseteq\mathcal F_\infty,
 \qquad \sigma(\mathcal A)=\sigma(\mathcal C)=\mathcal F_\infty.
\]
Consistency on bounded rectangles implies, for every Borel \(B\subseteq
\mathbb R^n\) and \(m\geq n\),
\[
                  \mu_m(B\times\mathbb R^{m-n})=\mu_n(B).
\]
Consequently
\[
 p_0(\pi_n^{-1}(B))\defeq\mu_n(B)
\]
is well defined on \(\mathcal C\), is finitely additive there,
and has total mass one. It is monotone and finitely subadditive.
Countable additivity is a further conclusion, proved below.
\end{hblemma}

\section{The compact approximation and diagonal step}
\label{sec:p15-cylinder-compactness}

\begin{hbdefinition}{def:p15-compact-approximation}{Compact sets and tight finite-dimensional probabilities}
A subset of Euclidean space is \emph{compact} when every cover by
open sets has a finite subcover. A probability measure \(\mu\) on
\(\mathbb R^d\) is \emph{tight} if for every \(\varepsilon>0\)
there is compact \(K\) with \(\mu(K)>1-\varepsilon\).
Closed and open below refer to the usual Euclidean topology.
\end{hbdefinition}

\begin{hblemma}{thm:p15-cylinder-compact-approximation}{proved=annex:p15.cylinder-compact-approximation}{The exact compact approximation needed for cylinders}
Every closed bounded Euclidean rectangle is compact. Finite unions
of such rectangles are compact and closed, and closed subsets of
compact sets are compact. Every Borel probability on \(\mathbb R^d\)
is tight.

If \(B\subseteq\mathbb R^d\) is a finite union of half-open
rectangles with the allowed infinite endpoints, and \(\mu\) is
a Borel probability, then for every \(\varepsilon>0\) there is
a finite union \(K\subseteq B\) of closed bounded rectangles
such that \(\mu(B\setminus K)<\varepsilon\). The empty union
is allowed. This is the compact inner approximation used in the
source's cylinder-algebra proof.
\end{hblemma}

\begin{hblemma}{thm:p15-cylinder-continuity}{proved=annex:p15.cylinder-continuity}{Continuity at the empty intersection}
For consistent \(\mu_n\) and the cylinder mass \(p_0\), if
\(B_n\in\mathcal A\) decrease and \(\bigcap_n B_n=\varnothing\),
then \(p_0(B_n)\downarrow0\).

The proof uses compact sets in finite-dimensional coordinate spaces
and a diagonal subsequence. A cylinder with a compact finite-dimensional
base still has unrestricted later coordinates; it is not being called
a compact subset of the full sequence space.
\end{hblemma}

\begin{hbtheorem}{thm:p15-kolmogorov-countable}{proved=annex:p15.kolmogorov-countable}{Kolmogorov extension for countably many real coordinates}
For every consistent family \((\mu_n)_{n\geq1}\) there exists
a unique probability measure \(P\) on
\((\Omega_\infty,\mathcal F_\infty)\) such that
\[
 P(\pi_n^{-1}(B))=\mu_n(B)
       \qquad(n\geq1,\ B\in\mathcal B(\mathbb R^n)).
\]
In particular its values on every rectangular cylinder are the values
specified in Durrett's Theorems 2.1.21 and A.3.1.
Every coordinate \(X_i\) is measurable, and every finite set of
coordinates has the law obtained by projecting \(\mu_m\), where
\(m\) is at least the largest index in that set. The measure is
on the displayed product sigma-algebra, without a silently added completion.
\end{hbtheorem}

\begin{hbcorollary}{thm:p15-independent-sequence-existence}{proved=annex:p15.independent-sequence-existence}{Independent sequences with prescribed real laws}
For any sequence of Borel probability laws \((\eta_i)_{i\geq1}\)
on \(\mathbb R\), there is a probability space carrying mutually
independent real variables \(X_i\) with those laws. One canonical
choice is the coordinate space above with finite-dimensional laws
\(\mu_n=\eta_1\times\cdots\times\eta_n\). This also realizes
every sequence of admissible real cumulative distribution functions
using their constructed Stieltjes measures.

For Borel sets \(E_i\subseteq\mathbb R\), the resulting measure satisfies
\[
 P(X_i\in E_i\text{ for every }i\geq1)
                 =\lim_{n\to\infty}\prod_{i=1}^n\eta_i(E_i).
\]
The limit is a decreasing limit in \([0,1]\). For fair binary
\(\eta_i=(\delta_0+\delta_1)/2\), this is exactly the law on
\(\mathbb R^{\mathbb N_+}\) of the dyadic sequence constructed in
Theorem~\ref{thm:p15-dyadic-independent-coins}, by uniqueness of the
finite-dimensional extension.
\end{hbcorollary}

\section{The source's binary coding, with its image and inverse}
\label{sec:p15-sequence-borel-coding}

\begin{hbdefinition}{def:p15-measurable-real-coding}{Metric and measurable coding conventions}
A metric \(\rho:M\times M\to[0,\infty)\) is zero exactly for equal points,
is symmetric, and satisfies the triangle inequality. Its open balls
\(\{y:\rho(x,y)<r\}\), for \(r>0\), generate its topology,
and \(\mathcal B(M)\) is the sigma-algebra
generated by its open sets. A metric space is \emph{complete} when
every Cauchy sequence converges in it; a sequence is Cauchy when for
every positive \(\varepsilon\), all sufficiently late pairwise
distances are less than \(\varepsilon\). Convergence means that distances
to the limit tend to zero. It is \emph{separable} when it has a countable
subset meeting every nonempty open ball. A homeomorphism is a bijection
continuous in both directions. For a subspace \(S\subseteq M\),
use the trace Borel sigma-algebra
\(\mathcal B(S)=\{S\cap B:B\in\mathcal B(M)\}\).

A \emph{measurable isomorphism} is a bijection with measurable forward
and inverse maps. A \emph{Borel real coding} of \((S,\mathcal S)\)
is such an isomorphism onto a Borel subset \(E\subseteq\mathbb R\)
with its trace Borel sigma-algebra. Spaces admitting such a coding are
called \emph{standard Borel spaces} here. The source calls the class
in its Theorem 2.1.22 \emph{nice}; the construction below supplies
both measurable directions and proves that its actual image is Borel.
\index{standard Borel space}\index{measurable isomorphism}
\end{hbdefinition}

\begin{hbdefinition}{def:p15-binary-sequence-code}{Canonical binary digits and diagonal interleaving}
Let \(Q=[0,1)^{\mathbb N_+}\) with its coordinate product
sigma-algebra. For \(t\in[0,1)\), define the canonical digit
\[
 b_k(t)=\lfloor2^kt\rfloor-2\lfloor2^{k-1}t\rfloor,
                         \qquad k\geq1.
\]
The convention chooses the expansion with infinitely many zeros,
including terminating zeros at dyadic numbers. Define a bijection of
pairs of positive indices with positive integers by
\[
 \theta(j,k)=\frac{(j+k-2)(j+k-1)}2+j.
\]
For \(x=(x_j)\in Q\), the source's diagonal binary code is
\[
 \phi_0(x)=\sum_{j,k\geq1}2^{-\theta(j,k)}b_k(x_j).
\]
The sum means the nonnegative series in increasing order of
\(\theta(j,k)\). Its first digit positions use
\((j,k)=(1,1),(1,2),(2,1),(1,3),(2,2),(3,1),\ldots\),
exactly the source's interleaving order. Put
\[
 E_0=\bigcap_{j,r\geq1}\bigcup_{k\geq r}
           \{t\in[0,1):b_{\theta(j,k)}(t)=0\}.
\]
Thus each decoded row is required to contain infinitely many zeros.
\end{hbdefinition}

\begin{hbtheorem}{thm:p15-sequence-borel-coding}{proved=annex:p15.sequence-borel-coding}{The exact binary image and a measurable inverse}
The product sigma-algebra on \(Q\) is the Borel sigma-algebra for
the metric
\(\delta(x,y)=\sum_{j\geq1}2^{-j}|x_j-y_j|\).
The map \(\phi_0\) is a measurable bijection from \(Q\) onto
the Borel set \(E_0\subseteq[0,1)\). Its inverse has coordinates
\[
 D_j(t)=\sum_{k\geq1}2^{-k}b_{\theta(j,k)}(t),
                  \qquad t\in E_0.
\]
Each \(D_j\) is measurable and takes values in \([0,1)\), and
\(D(t)=(D_j(t))\) is measurable into \(Q\). The point zero is
included; the point one is excluded; no ambiguous binary expansion
or unspecified exceptional input remains.
\end{hbtheorem}

\begin{hbtheorem}{thm:p15-metric-borel-representation}{proved=annex:p15.metric-borel-representation}{Borel subsets of complete separable metric spaces have real codings}
Let \((M,\rho)\) be a complete separable metric space and let
\(S\subseteq M\) be Borel, with its trace Borel sigma-algebra.
Then \(S\) has a Borel real coding. For nonempty \(S\), one
explicit construction is as follows. Choose a countable dense sequence
\((q_j)\) in \(S\), let \(H=\overline S\) be its closure in \(M\), and put
\[
 d(x,y)=\frac{\rho(x,y)}{1+\rho(x,y)},\qquad
 \psi(x)=(d(x,q_j))_{j\geq1}\quad(x\in H).
\]
The metric \(d\) is complete on \(H\), is less than one,
and gives the same topology as \(\rho\). The map \(\psi\)
is a homeomorphism of \(H\) onto the Borel subset
\[
 C=\left\{z\in Q:
 \begin{array}{l}
 |z_i-z_j|\leq d(q_i,q_j)\leq z_i+z_j\text{ for all }i,j,\\
 \inf_j z_j=0
 \end{array}\right\}.
\]
Write \(R:C\to H\) for its inverse. Then \(K=\psi(S)\)
is Borel in \(Q\), and
\[
 E=\{t\in E_0:D(t)\in K\},\qquad
 \phi=\phi_0\circ\psi|_S:S\longrightarrow E
\]
is the required coding, with inverse \(t\mapsto R(D(t))\).
For empty \(S\), take the empty map onto the empty Borel image.
This completes the source's Theorem 2.1.22 construction with the
image and inverse qualifications made explicit.
\end{hbtheorem}

\section{Transporting the extension to the specified measurable spaces}
\label{sec:p15-coded-space-extension}

\begin{hbtheorem}{thm:p15-coded-space-extension}{proved=annex:p15.coded-space-extension}{Countable extension through Borel real codings}
For each \(i\geq1\), let \((S_i,\mathcal S_i)\) admit a
Borel real coding \(\phi_i:S_i\to E_i\subseteq\mathbb R\).
Suppose probability laws \(\nu_n\) on the first \(n\) product
spaces are consistent under dropping the last coordinate, on all
measurable finite-product events. Then there is a unique probability
on \(\prod_{i\geq1}S_i\), equipped with its coordinate product
sigma-algebra, having those finite-dimensional laws.

In particular the construction applies when each \(S_i\) is a
Borel subset of a complete separable metric space. The proof first
constructs the real-coordinate measure, proves its concentration on
the measurable set \(\prod_i E_i\), and applies the measurable
coordinate inverses there. No extension assertion for arbitrary
measurable spaces is inferred from this theorem.
\end{hbtheorem}

\begin{hbexample}{def:p15-countable-product-boundaries}{Inconsistent specifications and a null infinite event}
If \(\mu_1=\delta_0\) but \(\mu_2\) is unit point mass at \((1,1)\), the
first-coordinate marginals disagree. No probability measure on real
sequences can have both laws: it would give probability one to each
of the disjoint events \(X_1=0\) and \(X_1=1\).

For the consistent fair-bit product, the event that the first \(n\)
bits are zero has probability \(2^{-n}\). Their decreasing
intersection, the event that every bit is zero, is nonempty but has
probability zero by continuity from above. Positive finite-cylinder
probabilities do not imply a positive probability for their infinite
intersection. This is an application of the constructed measure,
not an additional infinite-product assumption.
\end{hbexample}

\begin{proofboundary}{def:p15-countable-product-next}{Further representation and limiting results}
The countable real extension and the source's stated measurable-space
coding are now proved. The subsequent source discussion of the cardinality
of uncountable Borel sets requires its separate perfect-set argument;
it is not supplied by a bare injective coding.
\Cref{sec:p15-radon-nikodym,sec:p15-conditional-expectation,sec:p15-regular-conditional-laws}
develop density representation and conditioning, while
\cref{sec:p15-limsup-events-and-zero-one-laws} supplies Borel--Cantelli
and both zero-one laws. General completion, remaining independence
exercises and later weak and strong laws require further proofs.
No uncountable-index extension or measure on an unrestricted power set
is asserted here.
\end{proofboundary}
